Before the internship it was asked to follow the online lectures of the CS231n
course from Stanford University. Two of the three assignements were completed,
they covered many aspects of machine learning: Image classification, kNN,
Softmax, FC Neural Network, Fully Connected Networks, Batch Normalization,
Dropout, CNNs, Network Visualization, RNNs. 

The results of these two assignements will be presented here.

\subsubsection{Assignment 1}

The first assignement was mostly about implementing first elements of machine
learning all from scratch, no pytorch was used.

The Support Vector Machine (SVM) is one layer of a Neural Network without
activation function, it was asked to implement one from scratch and train it.
One exercice was to visualize the weights of the weight matrix after training a
SVM on the CIFAR-10 dataset with cross validation and softmax loss, it gives a
kind of average of the images of each class in \cref{fig:svm-visu}. The
cross-validation was used to choose the learning rate and regularization
strength hyperparameters, as shown in \cref{fig:cross-validation}.
We achieve almost 40 percent with the SVM.

\insertfigure
    {assets/images/svm-visu}
    {width=0.7\textwidth}
    {Visualization of the SVM's weights after training}
    {fig:svm-visu}

\insertfigure
    {assets/images/cross-validation}
    {width=0.7\textwidth}
    {Cross-validation to choose learning rate and regularization strength
    parameters}
    {fig:cross-validation}

Then, a two layers neural network was implemented, all from scratch, including
gradients calculus for backward propagation. ReLU activation functions were used
and our fully working two layers neural network achieved 52 percent accuracy on
the CIFAR-10 dataset. \cref{fig:twolayers-train} shows the training of it.
\cref{fig:two-layers-visu} shows the visualization of the weights of the first
layer, more abstract elements are learnt on this layer.

\insertfigure
    {assets/images/twolayers-train}
    {width=0.7\textwidth}
    {Training of the two layers model on CIFAR-10}
    {fig:twolayers-train}

\insertfigure
    {assets/images/two-layers-visu}
    {height=0.3\textheight}
    {Visualization of the weight of a trained two layers model on CIFAR-10}
    {fig:two-layers-visu}

One exercice on feature extraction asked to extract the color histograms of the
images and use that as an input to the classifier. This achieved a 58 percent
accuracy on the test set.

\subsubsection{Assignment 2}

In this assignement, more advanced layers were implemented from scratch before
moving to using PyTorch.

At first Batch Normalization is implemented from scratch (with forward pass and
backward pass), \cref{fig:batchnorm-efficiency} shows how it improves the
training.

\insertfigure
    {assets/images/batchnorm-efficiency}
    {height=0.3\textheight}
    {Batchnormalization improving training speed}
    {fig:batchnorm-efficiency}

Then, the dropouts layer was implemented from scratch to prevent overfitting,
\cref{fig:dropout} shows a less overfitting model with 0.25 dropout.

\insertfigure
    {assets/images/dropout}
    {height=0.3\textheight}
    {Batchnormalization improving training speed}
    {fig:dropout}

Finally, the last layer implemented from scratch is the convolution, spatial
batch normalization and spatial group normalization layers (forward pass and
backward pass included). The filters of the first layer of a CNN can be
visualized \cref{fig:visu-cnn}, they correspond to some edges/textures
detection.

\insertfigure
    {assets/images/visu-cnn}
    {height=0.3\textheight}
    {Visualization of the filters of the first layer of CNN trained on CIFAR-10}
    {fig:visu-cnn}

Finally, it was asked to use PyTorch, implementing a two-layers neural network
in the PyTorch framework. A final task was to try and implement the best
CIFAR-10 classifier. A 6 layers AlexNet like architecture showed 78 percent
accuracy. 